\section{Ordinary powers and relative flat flags}
\label{sec:supplement-relative-powers}
\begin{corollary}[Ordinary powers]\label{cor:powers}
For $n\ge1$, write $D^{[n]}=V(\II_D^n)$ and $E^{[n]}=V(\II_E^n)$.
Then
\begin{equation}\label{eq:powers}
 \II_D^n=\II_\Delta^n\II_E^n
           =\II_\Delta^n\cap\II_E^{2n}
\end{equation}
is an irredundant primary decomposition on $U$. With
$\mathcal T_n=\operatorname{tor}_{\eta_\Delta}(\OO_{D^{[n]}})$,
\begin{equation}\label{eq:power-sequence}
 0\longrightarrow\mathcal T_n=\OO_{E^{[n]}}(-n\Delta)
 \longrightarrow\OO_{D^{[n]}}\longrightarrow\OO_{n\Delta}
 \longrightarrow0.
\end{equation}
Its annihilator defines $E^{[n]}$. The generic multiplicity on
$\Delta$ is $n$, the nilradical index along $E$ is $2n$, and the
length of $\mathcal T_n$ at $\eta_E$ is $n(n+1)/2$.
\end{corollary}
\begin{proof}
In the local ring let $P=(t,f)$. The associated graded ring of $P$
is a polynomial ring in two variables over $\OO_E$, and the initial
form of $t$ is one of those variables. Hence
$(P^{2n}:t^n)=P^n$, proving
$(t^n)\cap P^{2n}=t^nP^n$. The same graded argument proves that these
powers are primary. The witnesses $t^n$ and $f^{2n}$ at $\eta_E$
prove irredundancy. Localizing at $(t)$ contracts the ideal to $(t^n)$,
so the generic torsion is $(t^n)/t^nP^n$, and cancellation gives its
annihilator $P^n$. The class of $t$ has index $2n$ at $E$, since
$t^{2n}$ belongs to the ideal but $t^{2n-1}$ does not. At the generic
point of $\Delta$, $f$ is a unit. Finally
$\OO_{U,\eta_E}/(t,f)^n$ has the $n(n+1)/2$ monomials of total
degree less than $n$ as successive graded basis elements.
\end{proof}

\begin{lemma}[Flat flags and arbitrary base change]\label{lem:flat-flag}
Let $U\to T$ be smooth, and let $E\subset\Delta\subset U$ be
relative smooth closed subschemes, each inclusion a Cartier divisor.
Put $L=\II_\Delta$ and $P=\II_E$. Then $\OO_U/P^n$ and
$\OO_U/L^n$ are flat over $T$. For $D_n=V(L^nP^n)$ the sequence
\begin{equation}\label{eq:relative-sequence}
 0\to L^n\otimes\OO_{E^{[n]}}\to\OO_{D_n}
        \to\OO_U/L^n\to0
\end{equation}
is exact and remains exact after any base change. Its base-changed
terms are the indicated ordinary powers and products. In particular
$D_n$ is flat over $T$.
\end{lemma}
\begin{proof}
The immersion $E\subset U$ is regular of codimension two and has
locally free conormal module. The conormal filtration of $\OO_U/P^n$
has quotients $\Sym^j(P/P^2)$ on $E$, for $0\le j<n$; these are
base-flat. Induction proves flatness. The analogous principal filtration
proves flatness of $\OO_U/L^n$. Multiplication by a local equation
$t^n$ of $L^n$ identifies the kernel in \eqref{eq:relative-sequence}
with $L^n\otimes\OO_{E^{[n]}}$, since $t$ is a non-zero-divisor.
Both outer terms are base-flat, so the middle term is base-flat.
For a base algebra $A\to A'$, the only possible left obstruction
in tensoring this short exact sequence is
$\Tor_1^A(\OO_U/L^n,A')$, which is zero. The sequences defining
$L$, $P$, and their powers stay exact for the same reason. Thus their
images after base change are the actual ideals generated by the
base-changed equations, and multiplication gives their stated products.
No assertion about specialization of selected primary components is
used.
\end{proof}

\begin{corollary}\label{cor:relative-flat}
Over $G_e^\circ$, the schemes $E$, $\Delta$, and $U$ form the smooth
relative flag of Lemma~\ref{lem:flat-flag}. All $D^{[n]}$ and
$E^{[n]}$ are flat, and \eqref{eq:power-sequence} commutes with
arbitrary base change. Each complex fibre has \eqref{eq:powers}.
\end{corollary}
\begin{proof}
Use the universal quotient $\mathcal S$ and the graph bundles
\eqref{eq:graph-delta}--\eqref{eq:graph-E} over the smooth family of
Proposition~\ref{prop:relative-open}. The Fitting elimination works
over the parameter ring, so it gives $t(t,f)$ relatively. Apply the
lemma, and apply Corollary~\ref{cor:powers} separately on each fibre
for the primary assertion.
\end{proof}
